
\section{Recordatorio sobre la teoría de stacks}
\label{XIII.0}

En lo que sigue utilizaremos la teoría de stacks expuesta en
\cite{XIII.1} y~\cite{XIII.2}. Nos limitamos al caso del sitio étale de un
esquema. Dado un esquema~$X$, denotemos por~$X_\et$ el sitio étale de~$X$.
Recordemos que un
\ifthenelse{\boolean{orig}}
{stack}
{stack $F$}
\index{stack|hyperpage}%
sobre~$X$ es una categoría fibrada sobre~$X_\et$ tal que, para todo
esquema~$X'$ étale sobre~$X$ y para todo par
\ifthenelse{\boolean{orig}}
{de objeto}
{de objetos}
$x$, $y$ de la fibra~$F_{X'}$, el prehaz~$\SheafHom_{X'}(x,y)$ sea un haz,
y tal que, para todo morfismo étale sobreyectivo $X''\to X'$, todo objeto
de~$F_{X''}$ provisto de un dato de descenso relativamente a $X''\to X'$
sea imagen inversa de un objeto de~$F_{X'}$.

Se denota por~$F(X')$ la categoría de las secciones cartesianas de~$F/X'$.
Más generalmente, si~$\Sch_X$ es la categoría de los esquemas sobre~$X$
provista de la topología étale, el stack~$F$ puede extenderse a un
stack~$\cal{F}$ sobre~$\Sch_X$ y, para todo morfismo $f\colon X'\to X$, se
denota también por~$F(X')$ la categoría de las secciones cartesianas de este
stack~$\cal{F}$ sobre~$X'$.

Un gerbe
\index{gerbe|hyperpage}%
es un stack tal que, para todo esquema~$X'$ étale sobre~$X$ y para todo par
de objetos $x$, $y$ de~$F_{X'}$, todo morfismo de~$x$ a~$y$ sea un
isomorfismo, que $x$ e~$y$ sean localmente isomorfos y que la familia de
objetos~$X'$ de~$X_\et$ tales que la fibra~$F_{X'}$ no sea vacía constituya
un recubrimiento de~$X_\et$. Por ejemplo, el stack de los torsores bajo un haz
de grupos es un gerbe que, además, posee una sección cartesiana.
Recíprocamente, un gerbe que posee una sección, \ie tal que existe un
objeto~$x$ de~$F_X$, es equivalente al stack de los torsores bajo el haz
\marginpar{346}
de grupos~$\SheafAut_X(x)$.

Existe una noción evidente de subgerbe y de subgerbe maximal de un
stack~$F$. Dada una sección cartesiana~$x$ de~$F(X)$, existe un único
subgerbe maximal~$G_x$ de~$F$ tal que~$x$ se factoriza a través de~$G_x$.
$G_x$ recibe el nombre de subgerbe generado por~$x$; por definición, es un
gerbe trivial.
\index{gerbe trivial|hyperpage}%
El prehaz~$SF$
\label{indnot:mb}\oldindexnot{$SF$ o $S(F)$|hyperpage}%
definido por
$$
SF(X') = \{ \text{subgerbes maximales de }F|X' \}
$$
es un haz, llamado haz de los subgerbes maximales de~$F$.
\index{haz de los subgerbes maximales|hyperpage}%
Sea~$O$ el prehaz definido por
$$
O(X') = \{ \text{clases de objetos de }F_{X'}\text{ salvo isomorfismo}
\} \text{.}
$$
Asociando a todo objeto~$x$ de~$F_{X'}$ el subgerbe maximal de~$F|X'$
generado por~$x$, se obtiene un morfismo
$$
O \to SF \text{;}
$$
según~\cite[III 2.1.4]{XIII.2}, este morfismo hace de~$SF$ un haz asociado
a~$O$.

Se dice que un stack~$F$ es \emph{constructible}
\index{constructible (stack)|hyperpage}%
(\emph{\resp ind-$\LL$-finito},
\index{stack ind-finito (\resp ind-$l$-finito)|hyperpage}%
donde~$\LL$ es un conjunto de números primos) si, para todo esquema~$X'$
étale sobre~$X$ y para todo objeto~$x$ de~$F_{X'}$, el haz
$\SheafAut_{X'}(x)$ posee la propiedad correspondiente
\cite[VII 2.2.1]{XIII.2}. Se dice que un stack es $1$-constructible
\index{stack $1$-constructible (\resp constructible)|hyperpage}%
si es constructible y si el haz de los subgerbes maximales es constructible.
